%
% new_tokens.tex
%
% A Z specification of the luxd-pty daemon's memory of retry tokens, the
% entry it keeps for each launch it admits from a `new` request's token.
% The intact design.
%
% Design record: DESIGN.md, DES-107 (token memory, "Formal model")
% Type-check:  fuzz -t docs/new_tokens.tex
% Model-check: scripts/verify_new_tokens_model.sh
%
\documentclass[a4paper,10pt,fleqn]{article}
\usepackage[margin=1in]{geometry}
\usepackage{fuzz}

\begin{document}

\title{The Daemon's Memory of Retry Tokens: a Z Specification}
\author{Formal model of the daemon's retry-token memory}
\date{November 2026}
\maketitle

\section{Introduction}

A caller of \texttt{new} retries after \texttt{timeout}, \texttt{provider\_lost}
or \texttt{fault} with the same token, because each of the three can follow a
launch that already happened.  The daemon therefore keeps one entry per launch
it admitted.  The property to establish is that a token that started a launch
never starts a second one while its entry is held, and that the entry is held
for exactly as long as the contract says: the 256 most recent admitted
launches, scaled here to $cap = 2$.

Launches are serialised in the daemon, so each operation below is one atomic
step.  The daemon adds the entry after every check that could refuse the launch
and before it forks.  Four outcomes follow.  The launch succeeds; the fork
fails (the entry is removed and the oldest entry the addition dropped is
restored); the fork succeeds but the session is never recorded (the entry stays
unbound, and the daemon logs token and process id); or the launch is refused
before the fork and adds nothing.  A repeated token with a session returns that
session; a repeated token without one is refused.

Three history variables belong to the specification, not to the daemon.
\texttt{recent} is the last $cap$ admitted launches, \texttt{nforks} counts the
forks of each token since the token left \texttt{recent} or the daemon
restarted, and the invariants compare the daemon's memory \texttt{held} with
them.  \texttt{held} gives each entry an age: the number of admitted launches
since it was added.

\section{Basic Types}

\begin{zed}
  [TOK] \\
  LAB ::= lbl1 | lbl2
\end{zed}

\begin{axdef}
  cap : \nat
\where
  cap = 2
\end{axdef}

\texttt{TOK} is a retry token.  \texttt{LAB} is a session label.  The
contract's 256 is scaled to $cap = 2$ with three tokens, which is enough to
evict one entry and to retry it.

\section{State}

\begin{schema}{Sys}
  held : TOK \pfun (0 \upto cap) \\
  bound, sv, kas : \power TOK \\
  lab : LAB \pfun TOK \\
  lv, surv, nforks : TOK \fun (0 \upto 2) \\
  recent : \seq TOK
\where
  bound \subseteq \dom held \\
  \# recent \leq cap \\
  \forall t : TOK @ \# (lab \rres \{ t \}) + surv~t \leq lv~t
\end{schema}

\texttt{held} is the memory: a token maps to the age of its entry.
\texttt{bound} are the held tokens whose session was recorded.  \texttt{lab}
maps a label to the token of the live session that owns it.  \texttt{lv} counts
the live processes of a token, and \texttt{surv} counts those that are
survivors of a restart: a survivor is a process, not a token, so it persists
through later clean restarts only while it is itself alive.  Labelled sessions
and survivors are distinct live processes, which the last predicate states.
\texttt{sv} (every token that has had a survivor), \texttt{kas}, \texttt{recent} and
\texttt{nforks} are history variables.

\begin{schema}{Init}
  Sys'
\where
  held' = \emptyset \\
  bound' = \emptyset \\
  sv' = \emptyset \\
  kas' = \emptyset \\
  lab' = \emptyset \\
  lv' = TOK \cross \{ 0 \} \\
  surv' = TOK \cross \{ 0 \} \\
  nforks' = TOK \cross \{ 0 \} \\
  recent' = \langle \rangle
\end{schema}

\section{Operations}

\begin{schema}{RepeatBound}
  \Delta Sys \\
  t? : TOK \\ l? : LAB
\where
  t? \in \dom held \\
  t? \in bound \\
  held' = held \\
  bound' = bound \\
  lab' = lab \\
  lv' = lv \\
  surv' = surv \\
  sv' = sv \\
  kas' = kas \\
  recent' = recent \\
  nforks' = nforks
\end{schema}

A repeated token whose session was recorded is answered with the original session, running or exited. Nothing changes.

\begin{schema}{RepeatUnrecorded}
  \Delta Sys \\
  t? : TOK \\ l? : LAB
\where
  t? \in \dom held \\
  t? \notin bound \\
  held' = held \\
  bound' = bound \\
  lab' = lab \\
  lv' = lv \\
  surv' = surv \\
  sv' = sv \\
  kas' = kas \\
  recent' = recent \\
  nforks' = nforks
\end{schema}

A repeated token whose process was forked but whose session was never recorded is refused as \texttt{rejected}. Nothing changes.

\begin{schema}{RefuseLabel}
  \Delta Sys \\
  t? : TOK \\ l? : LAB
\where
  t? \notin \dom held \\
  l? \in \dom lab \\
  held' = held \\
  bound' = bound \\
  lab' = lab \\
  lv' = lv \\
  surv' = surv \\
  sv' = sv \\
  kas' = kas \\
  recent' = recent \\
  nforks' = nforks
\end{schema}

A launch refused because its label is held by a live session. The refusal precedes the fork and adds no entry.

\begin{schema}{RefuseConfig}
  \Delta Sys \\
  t? : TOK \\ l? : LAB
\where
  t? \notin \dom held \\
  held' = held \\
  bound' = bound \\
  lab' = lab \\
  lv' = lv \\
  surv' = surv \\
  sv' = sv \\
  kas' = kas \\
  recent' = recent \\
  nforks' = nforks
\end{schema}

Any other refusal before the fork: runtime launchers off, unknown launcher. No entry is added.

\begin{schema}{LaunchOk}
  \Delta Sys \\
  t? : TOK \\ l? : LAB
\where
  t? \notin \dom held \\
  l? \notin \dom lab \\
  lv~t? < 2 \\
  nforks~t? < 2 \\
  held' = \{ x : TOK; a : 0 \upto cap | (x \mapsto a) \in (\{ t? \} \ndres held) \land a + 1 < cap @ x \mapsto a + 1 \} \cup \{ t? \mapsto 0 \} \\
  bound' = (bound \cap \dom held') \cup \{ t? \} \\
  lab' = lab \cup \{ l? \mapsto t? \} \\
  lv' = lv \oplus \{ t? \mapsto lv~t? + 1 \} \\
  surv' = surv \\
  sv' = sv \\
  kas' = kas \\
  ((\# recent < cap \land recent' = recent \cat \langle t? \rangle) \lor (\# recent = cap \land (\exists h : TOK; r : \seq TOK @ recent = \langle h \rangle \cat r \land recent' = r \cat \langle t? \rangle))) \\
  nforks' = (\ran recent' \dres (nforks \oplus \{ t? \mapsto nforks~t? + 1 \})) \cup ((TOK \setminus \ran recent') \cross \{ 0 \})
\end{schema}

An admitted launch whose fork and session both succeed. The entry is added before the fork and bound to the session.

\begin{schema}{LaunchForkFails}
  \Delta Sys \\
  t? : TOK \\ l? : LAB
\where
  t? \notin \dom held \\
  l? \notin \dom lab \\
  held' = held \\
  bound' = bound \\
  lab' = lab \\
  lv' = lv \\
  surv' = surv \\
  sv' = sv \\
  kas' = kas \\
  recent' = recent \\
  nforks' = nforks
\end{schema}

The fork fails. The entry added before the fork is removed and the oldest entry its addition dropped is restored, so the memory is as it was, and the reply is \texttt{fault}.

\begin{schema}{LaunchUnrecorded}
  \Delta Sys \\
  t? : TOK \\ l? : LAB
\where
  t? \notin \dom held \\
  l? \notin \dom lab \\
  lv~t? < 2 \\
  nforks~t? < 2 \\
  held' = \{ x : TOK; a : 0 \upto cap | (x \mapsto a) \in (\{ t? \} \ndres held) \land a + 1 < cap @ x \mapsto a + 1 \} \cup \{ t? \mapsto 0 \} \\
  bound' = bound \cap \dom held' \\
  lab' = lab \\
  lv' = lv \oplus \{ t? \mapsto lv~t? + 1 \} \\
  surv' = surv \\
  sv' = sv \\
  kas' = kas \\
  ((\# recent < cap \land recent' = recent \cat \langle t? \rangle) \lor (\# recent = cap \land (\exists h : TOK; r : \seq TOK @ recent = \langle h \rangle \cat r \land recent' = r \cat \langle t? \rangle))) \\
  nforks' = (\ran recent' \dres (nforks \oplus \{ t? \mapsto nforks~t? + 1 \})) \cup ((TOK \setminus \ran recent') \cross \{ 0 \})
\end{schema}

The entry is added, the process forks, and the session is never recorded.

\begin{schema}{ExitLabelled}
  \Delta Sys \\
  l? : LAB
\where
  l? \in \dom lab \\
  lv~(lab~l?) > 0 \\
  held' = held \\
  bound' = bound \\
  lab' = \{ l? \} \ndres lab \\
  lv' = lv \oplus \{ lab~l? \mapsto lv~(lab~l?) - 1 \} \\
  surv' = surv \\
  sv' = sv \\
  kas' = kas \\
  recent' = recent \\
  nforks' = nforks
\end{schema}

A recorded session exits and frees its label. Its token stays held and bound.

\begin{schema}{ExitStray}
  \Delta Sys \\
  t? : TOK
\where
  lv~t? - surv~t? > \# (lab \rres \{ t? \}) \\
  held' = held \\
  bound' = bound \\
  lab' = lab \\
  lv' = lv \oplus \{ t? \mapsto lv~t? - 1 \} \\
  surv' = surv \\
  sv' = sv \\
  kas' = kas \\
  recent' = recent \\
  nforks' = nforks
\end{schema}

A process with no recorded session, one that did not survive a restart, exits.

\begin{schema}{ExitSurvivor}
  \Delta Sys \\
  t? : TOK
\where
  surv~t? > 0 \\
  held' = held \\
  bound' = bound \\
  lab' = lab \\
  lv' = lv \oplus \{ t? \mapsto lv~t? - 1 \} \\
  surv' = surv \oplus \{ t? \mapsto surv~t? - 1 \} \\
  sv' = sv \\
  kas' = kas \\
  recent' = recent \\
  nforks' = nforks
\end{schema}

A shell that survived a restart exits. It no longer protects any later shell from a clean restart.

\begin{schema}{RestartClean}
  \Delta Sys
\where
  held' = \emptyset \\
  bound' = \emptyset \\
  lab' = \emptyset \\
  lv' = surv \\
  surv' = surv \\
  sv' = sv \\
  kas' = kas \cup \{ x : TOK | x \in sv \land surv~x = 0 \land lv~x > 0 \} \\
  recent' = \langle \rangle \\
  nforks' = TOK \cross \{ 0 \}
\end{schema}

The daemon restarts. The memory is cleared and every shell the restart can end is ended: each live process that is not itself a survivor of an earlier restart. A survivor is out of its reach and persists. \texttt{kas} records the tokens whose earlier survivor had gone and whose later shell this restart ended.

\begin{schema}{RestartSurvivor}
  \Delta Sys \\
  s? : \power TOK
\where
  s? \neq \emptyset \\
  \forall x : s? @ lv~x - surv~x > 0 \land surv~x < 2 \\
  held' = \emptyset \\
  bound' = \emptyset \\
  lab' = \emptyset \\
  lv' = surv' \\
  surv' = surv \oplus \{ x : s? @ x \mapsto surv~x + 1 \} \\
  sv' = sv \cup s? \\
  kas' = kas \cup \{ x : TOK | x \in sv \land surv~x = 0 \land lv~x > 0 \} \\
  recent' = \langle \rangle \\
  nforks' = TOK \cross \{ 0 \}
\end{schema}

The daemon restarts and a shell for each token in $s?$ survives SIGKILL. The daemon logs its token and process id and exits with a failure status; the survivor has no history.

\section{Invariants}

Each invariant is checked as a violation predicate in
\texttt{scripts/verify\_new\_tokens\_model.sh}.

\begin{description}
  \item[T1, a held token never yields a second launch.]
    $\forall t : TOK @ nforks~t \leq 1$.
  \item[T2, refusals and repeats add no entry.]
    $\dom held \subseteq \ran recent$: every entry belongs to an admitted
    launch of the last $cap$.
  \item[T3, the oldest is dropped only at the cap.]
    $\{ t : \ran recent | nforks~t > 0 \} \subseteq \dom held$: no token that
    forked in the last $cap$ admitted launches has been dropped early.  With T2
    and $\# recent \leq cap$ the memory is exactly the entries of the last
    $cap$ launches that forked.
  \item[T4, a launch with no process can be retried.]
    $\forall t : \dom held @ nforks~t > 0$: an entry exists only for a token
    whose launch forked.
  \item[T5, the memory is bounded.] $\# (\dom held) \leq cap$.
\end{description}

\section{Reachability}

An invariant that no state could violate proves nothing, so
\texttt{verify\_new\_tokens\_model.sh} also asks ProB to reach: a live shell
whose entry was evicted; a forked token with no recorded session; an exited
session whose token is still held; a full memory with a failed fork pending;
a survivor and its retry both live (the accepted residual case); and each
operation enabled at least once.

\end{document}
